% Part: lambda-calculus
% Chapter: lambda-definability
% Section: truth-values

\documentclass[../../../include/open-logic-section]{subfiles}

\begin{document}

\olfileid[pt]{lam}{ldf}{tvr}
\olsection{Valores de verdade e relações}

Podemos codificar valores de verdade no cálculo lambda puro da seguinte maneira:
\begin{align*}
  \fn{true} & \ident \lambd[x][\lambd[y][x]]\\
  \fn{false} & \ident \lambd[x][\lambd[y][y]]
\end{align*}

Os valores de verdade são representados como \emph{seletores}, isto é, funções que
recebem dois argumentos e devolvem um deles. O valor de verdade
$\fn{true}$ seleciona seu primeiro argumento, e $\fn{false}$, o
segundo. Por exemplo, $\fn{true}\, M N$ sempre se reduz a $M$, ao passo que
$\fn{false}\, M N$ sempre se reduz a~$N$.

\begin{defn}
Dizemos que uma relação $R \subseteq \Nat^k$ é !!{lambda definable} se existe
um termo~$R$ tal que
\begin{align*}
  R\, \num{n_1} \dots \num{n_k} & \bred \fn{true}
  \intertext{sempre que $R(n_1, \dots, n_k)$, e}
  R\, \num{n_1} \dots \num{n_k} & \bred \fn{false}
\end{align*}
caso contrário.
\end{defn}

Por exemplo, a relação $\fn{IsZero} = \{0\}$, que vale para $0$
e somente para $0$, é !!{lambda definable} por
\[
  \fn{IsZero} \ident \lambd[n][n (\lambd[x][\fn{false}])\, \fn{true}].
\]
Como isso funciona? Como os numerais de Church são definidos como iteradores
(funções que aplicam seu primeiro argumento $n$~vezes ao segundo),
fixamos o valor inicial como $\fn{true}$ e, a cada passo da
iteração, devolvemos $\fn{false}$ independentemente do resultado da iteração
anterior. Esse passo será aplicado ao valor inicial $n$
vezes, e o resultado será $\fn{true}$ se, e somente se, o passo não for
aplicado nenhuma vez, isto é, quando $n = 0$.

Com base nessa representação dos valores de verdade, podemos ainda
definir algumas funções de verdade. Eis duas delas, as representações da
negação e da conjunção:
\begin{align*}
  \fn{Not} & \ident \lambd[x][x\, \fn{false}\, \fn{true}]\\
  \fn{And} & \ident \lambd[x][\lambd[y][xy \,\fn{false}]]
\end{align*}
A função ``$\fn{Not}$'' recebe um argumento e devolve
$\fn{true}$ se o argumento for $\fn{false}$, e $\fn{false}$ se o
argumento for~$\fn{true}$. A função ``$\fn{And}$'' recebe dois
valores de verdade como argumentos e deve devolver $\fn{true}$ se, e somente se, ambos os
argumentos forem~$\fn{true}$. Os valores de verdade são representados como seletores
(descritos acima); portanto, quando $x$ é um valor de verdade e é aplicado a dois
argumentos, o resultado será o primeiro argumento se $x$ for $\fn{true}$,
e o segundo argumento caso contrário. Agora, $\fn{And}$ recebe seus dois
argumentos $x$ e $y$ e passa, em resposta, $y$ e $\fn{false}$ a
seu primeiro argumento~$x$. Supondo que $x$ seja um valor de verdade, o resultado
será avaliado como~$y$ se $x$ for $\fn{true}$, e como $\fn{false}$ se $x$
for $\fn{false}$, que é precisamente o que se deseja.

Observe que supomos aqui que apenas valores de verdade são usados como argumentos
de $\fn{And}$. Se outros termos lhe forem passados, o resultado (isto é, a
forma normal, se existir) pode muito bem não ser um valor de verdade.

\begin{prob}
  Defina as funções $\fn{Or}$ e $\fn{Xor}$ que representam as funções de verdade
  da disjunção inclusiva e da disjunção exclusiva usando a codificação
  dos valores de verdade como termos-$\lambd$.
\end{prob}

\end{document}
